%%%%%%%%%%%%%%%%%%%%%%%%%%%%%%%%%%%%%%%%%%%%%%%%%%%%%%%%%%%%%%%%%%%%%%%%%%%%%%
%	Plantilla para Presentaciones, UTFSM, Chile
%	=============================================
%
% Autor:
%       Jaime C. Rubin-de-Celis <jaime@rubin-de-celis.com>
%
% Fecha:
%       $Date: 2022-03-28 15:23:54 -0300 (28 Marzo 2022) $
%
% Versión:
%       1.0
%
% Licencia:
%       Copyright (c) 2018-2022, Jaime C. Rubin-de-Celis
%       The MIT License
%
% Nota:
%       Las imágenes son propiedad intelectual de la UTFSM.
%
% Uso:
%       Ver archivo adjunto README.md.
%       Si todo lo demás falla, recuerde "RTFM".
%
%%%%%%%%%%%%%%%%%%%%%%%%%%%%%%%%%%%%%%%%%%%%%%%%%%%%%%%%%%%%%%%%%%%%%%%%%%%%%%

\documentclass[
	9pt,
    aspectratio=169,
    professionalfonts,
	 handout % Uncomment to produce Handouts
]{beamer}

\usepackage[utf8,utf8x]{inputenc}
%\usepackage[scale=2]{ccicons} % Icons for CC-BY-SA
\usepackage{pgf-pie}

%%% HANDOUTS
% Uncomment the following lines to prepare '3 slides in 1' Handouts
% It is also advisable to add the option [handout] to the document class.
% Options: `2 slides in 1`, `3 slides in 1` `4 slides in 1`
%\usepackage{handoutWithNotes}
%\pgfpagesuselayout{3 on 1 with notes}[letterpaper,border shrink=5mm]


%%% CONFIGURATION
\newcommand{\TheLogo}{figures/logoutfsmind_v.png}	%Required
\newcommand{\TheBackground}{backgrounds/usmCCFrontis_transparent.jpg}
%\newcommand{\TheBackground}{backgrounds/bgWhiteAbstract.jpg}

\title{Título de la Presentación}
\subtitle{Subtítulo de la Presentación}
\author{Nombre del Autor}
\institute{Universidad Técnica Federico Santa María}
\date{\today}
\newcommand{\authoremail}{nombre.apellido@usm.cl}

%%% Transitions
\mode<presentation>
{
	\usetheme{UTFSM}
	\setbeamercovered{transparent = 28}
	\beamerdefaultoverlayspecification{<+->}
	% Global Background
	\usebackgroundtemplate{%
    	\includegraphics[width=\paperwidth,height=\paperheight]{\TheBackground}%
	}
}

\DeclareMathOperator*{\maxim}{Max}
\DeclareMathOperator*{\argmax}{arg\,max}
\numberwithin{equation}{section}



%---------------------------------------------------------------------------
%%% PDF OPTIONS (HYPERREFERENCES)
\hypersetup{
		colorlinks=false
		,linkcolor=black!80
		,urlcolor=blue
		,citecolor=red
		,breaklinks=true
  		,pdfinfo={  
    Title={Plantillas Latex UTFSM},  
    Subject={Aprende LaTeX a través del Ejemplo},
    Author={Jaime C. Rubin-de-Celis}
    Keywords={Beamer Template},
		Producer={JCR's LaTeX, http://www.sglms.com/},
		pdfpagemode=FullScreen,
		pdfmenubar=false,
		pdftoolbar=false
  }  
}
%---------------------------------------------------------------------------


%---------------------------------------------------------------------------
%%% FIGURES, GRAPHICS AND COLORS
\usepackage{wrapfig}			% flooating Figures and Tables
	% Colors
	\definecolor{gray80}{gray}{0.80}
	\definecolor{gray40}{gray}{0.40}
%---------------------------------------------------------------------------

%---------------------------------------------------------------------------
%%% TIKZ
\usepackage{tikz}				% to hard to explain ;-)
\usetikzlibrary{arrows,positioning,fit,shapes,calc,backgrounds}
\usetikzlibrary{decorations.pathreplacing}
\tikzstyle{line} = [draw, thick, gray40, ->, >=latex]
% ATTENTION: library conflicts with pgfpages!!!
%---------------------------------------------------------------------------


%---------------------------------------------------------------------------
%%% Begin Document
%---------------------------------------------------------------------------

\begin{document}

\frame{\titlepage}

%!TEX root = beamer  .tex

%==============================================
\section{Introduction}
%==============================================
\begin{frame}{\insertsectionhead}
  \framesubtitle{A short introduction to Trigon}
  \themename is a modern, elegant and versatile theme for Beamer, following UTFSM's
  corporate image guidelines and colors
  (\href{https://www.usm.cl}{\textsc{UTFSM}}). 
  \vfill
  \themename comes with lots of nice extra features
  \begin{itemize}
    \item Compliance with UT
    \item Simple customization of theme colors
    \item Lots of convenient options to tweak the design
  \end{itemize}
\end{frame}


%==============================================
\section{Layout}
%==============================================

\subsection{Layout variations}

\begin{frame}[fragile=singleslide]{\insertsectionhead}
  \framesubtitle{\insertsubsectionhead}
  The general style for the title, section and regular frames can be changed
  easily with simple options. Here are some examples for the title page

\begin{figure}
\centering     %%% not \center
\subfigure[Figure A]{\label{fig:b}\includegraphics[width=50mm]{figures/logoutfsmind_h.png}}
\subfigure[Figure B]{\label{fig:a}\includegraphics[width=30mm]{figures/logoutfsmind_v.png}}

\subfigure[Figure C]{\label{fig:a}\includegraphics[width=30mm]{figures/thumbnail.png}}
\caption{UTFSM Logos}
\end{figure}

\end{frame}

%--------------------------------------
\subsection{Fonts}

\begin{frame}
  \frametitle{\insertsectionhead}
  \framesubtitle{\insertsubsectionhead}
  This theme is using (Linux) Libertine fonts for all elements. This can be changed in theme template.
  \vfill
  Emphasis can be added by using \textbf{bold} typeface, \textit{italic},
  \alert{alert} or {\color{USMBlue}{simple colors}}.
  \vfill
  Equations are typsetted with this font as well
  \begin{equation*}
    F(x|\mu,s) = \int_{-\infty}^x s^{-1}\left(1+e^{-\frac{v-\mu}{s}}\right)^{-2} e^{-\frac{v-\mu}{s}}\;\mathsf{d}v = \frac{1}{1+e^{-\frac{x-\mu}{s}}}
  \end{equation*}
\end{frame}


%==============================================
\section{Elements}
%==============================================
\subsection{Charts}
\begin{frame}{\insertsectionhead}
	\framesubtitle{\insertsubsectionhead}
	\begin{columns}[c, onlytextwidth]
		\column{0.47\textwidth}
		Use the theme color \texttt{USMBlue}, \texttt{USMRed}, \texttt{USMYellow} and
		\texttt{USMYellow} to have charts directly fit the main theme of presentation.
		\vfill
		\begin{itemize}
			\item Easy variants using \texttt{color!x} to lighten or darken the colors
		\end{itemize}
		\hfill
		\column{0.47\textwidth}
		\center
		\resizebox{0.9\textwidth}{!}{%
			\begin{tikzpicture}
				\pie[color={USMBlue, USMRed!75, USMYellow!50, USMRed!25},
				rotate=90, hide number, text= legend]
				{50/USMBlue, 32/USMRed!75, 16/USMYellow!50, 2/USMRed!25}
			\end{tikzpicture}
		}
	\end{columns}
\end{frame}

\subsection{Lists}
\begin{frame}{\insertsectionhead}
  \framesubtitle{\insertsubsectionhead}
  \begin{columns}[T,onlytextwidth]
    \column{0.28\textwidth}
    Items
    \begin{itemize}
      \item Item 1
        \begin{itemize}
          \item Subitem 1
          \item Subitem 2
        \end{itemize}
      \item Item 2
      \item Item 3
    \end{itemize}

    \column{0.42\textwidth}
    Enumerations
    \begin{enumerate}
      \item The Fellowship of the Ring,
      \item The Two Towers,
      \item The Return of the King.
    \end{enumerate}

    \column{0.30\textwidth}
    Descriptions
    \begin{description}
      \item[UTFSM] Modern. \item[Default] Outdated.
    \end{description}
  \end{columns}
\end{frame}



%--------------------------------------
\subsection{Figures}
\begin{frame}
  \frametitle{\insertsectionhead}
  \framesubtitle{\insertsubsectionhead}

  \begin{figure}
    \newcounter{density}
    \setcounter{density}{20}
    \begin{tikzpicture}
      \def\couleur{USMYellow}
      \def\color{USMRed}
      \path[coordinate] (0,0)  coordinate(A)
        ++( 60:5cm) coordinate(B)
        ++(-60:5cm) coordinate(C);
      \path[coordinate] (0,0)  coordinate(D)
        ++(60:5cm) coordinate(E)
        ++(180:5cm) coordinate(F);
      \draw[fill=\couleur!\thedensity] (A) -- (B) -- (C) -- cycle;
      \draw[fill=\color!\thedensity] (D) -- (E) -- (F) -- cycle;
      \foreach \x in {1,...,15}{%
        \pgfmathsetcounter{density}{\thedensity+10}
        \setcounter{density}{\thedensity}
        \path[coordinate] coordinate(X) at (A){};
        \path[coordinate] (A) -- (B) coordinate[pos=.15](A)
          -- (C) coordinate[pos=.15](B)
          -- (X) coordinate[pos=.15](C);
        \draw[fill=\couleur!\thedensity] (A)--(B)--(C)--cycle;
      }
      \setcounter{density}{20}
      \foreach \x in {1,...,15}{%
        \pgfmathsetcounter{density}{\thedensity+10}
        \setcounter{density}{\thedensity}
        \path[coordinate] coordinate(X) at (D){};
        \path[coordinate] (D) -- (E) coordinate[pos=.15](D)
          -- (F) coordinate[pos=.15](E)
          -- (X) coordinate[pos=.15](F);
        \draw[fill=\color!\thedensity] (D)--(E)--(F)--cycle;
      }
    \end{tikzpicture}
    \caption{Rotated triangles from
    \href{http://www.texample.net/tikz/examples/rotated-triangle/}{texample.net}.}
  \end{figure}
\end{frame}

%--------------------------------------
\subsection{Tables}
\begin{frame}
  \frametitle{\insertsectionhead}
  \framesubtitle{\insertsubsectionhead}
  \begin{table}[H]
    \centering
    \caption{A simple table example}
    \begin{tabular}{@{} lccc @{}}
      \toprule
      & \textbf{Velocity} & \textbf{Angle}  & \textbf{Vertical force} \\
      & $U$ & $\alpha$  & $F_z$ \\
      & [m/s] & [$^\circ$]  & [N] \\
      \midrule
      2D simulation  & 9 & 2 & 9.23 \\
      3D simulation  & 10.0 & 3 & 15.039 \\
      Experiment A   & 11.31 & 2.5 & 13.2 \\
      Experiment B   & 11.26 & 2.7 & 12.6 \\
      Experiment C   & 11.33 & 2.47 & 13.6 \\
      \bottomrule
    \end{tabular}
  \end{table}

\end{frame}

%--------------------------------------
\subsection{Blocks}
\begin{frame}
  \frametitle{\insertsectionhead}
  \framesubtitle{\insertsubsectionhead}
  \begin{block}{Regular block}
    Just a regular block
  \end{block}
  \begin{alertblock}{Alert block}
    Some important thing
  \end{alertblock}
  \begin{exampleblock}{Example block}
    No difference with regular block to avoid excessive distraction
  \end{exampleblock}
\end{frame}

%--------------------------------------
\subsection{Frame footer}

%
\begin{frame}
    \frametitle{\insertsectionhead}
    \framesubtitle{\insertsubsectionhead}
	Footnote example\footnote{This is a footnote.}.
	\begin{example}
		Ejemplo\footnote{This is a footnote.} ...
		\begin{equation*}
			E =mc^2
		\end{equation*}
	\end{example}
\end{frame}


\begin{frame}{References}
  Some references to showcase 
  \citet{Knuth92},
  \citet{Simpson}
  \citet{greenwade93}
  \citet{ConcreteMath}
\end{frame}

%==============================================
\section{Conclusion}
%==============================================
\begin{frame}{Summary}

  Get the source of this theme and the demo presentation from

  \begin{center}\url{https://github.com/jaimercz/utfsm-beamer}\end{center}

  \themename is licensed under a
  \href{https://opensource.org/licenses/MIT}{MIT LICENCE}. Images are intellectual
  property of UTFSM \href{www.usm.cl}{USM}.

\end{frame}

%==============================================
%\begin{frame}[standout]
%  Questions?
%\end{frame}

\appendix

\begin{frame}[fragile]{Backup slides}
  Sometimes, it is useful to add slides at the end of your presentation to
  refer to during audience questions.

  The best way to do this is to include the \verb|appendixnumberbeamer|
  package in your preamble and call \verb|\appendix| before your backup slides.

  \themename will automatically turn off slide numbering and progress bars for
  slides in the appendix.
\end{frame}

% The following line is very important! There is a bug with hyperref and overlays.
% UTFSM theme uses <+-> specifications, but hyperref gets confused with this. 
% So we need to reset the specification to make it work.
\beamerdefaultoverlayspecification{<*>}
\begin{frame}[allowframebreaks]{References}

   \bibliography{biblio}

\end{frame}



\end{document}
